\scanpage{113}
\mspage{56}

Le lien avec la \deleted{\unsure{definition|définition}}\inserted{description} (11) est évidemment que, si $\unsure{\zeta|\xi}_X$ provient d'un $\unsure{p^t|pt}$ géom $\unsure{\zeta|\xi}_T$ de $T$, alors on aura des iso can

\[
\begin{aligned}
(20)\quad &\mathcal E_{X/Y}^{\unsure{\zeta|\xi}_T}\simeq\mathcal E_{X/Y}^T\\[0.5ex]
&\quad\rotatebox{90}{$\simeq$}\\[0.5ex]
&\Pi_1(X_{\unsure{\zeta|\xi}_Y},\unsure{\zeta|\xi}_T)
\end{aligned}
\]

\deleted{\unsure{donc}} --- les $\Pi_1(X_{\unsure{\zeta|\xi}_Y},\unsure{\zeta|\xi}_T)$ des fibres géom, relatifs aux divers pts géom de $\mathcal E_{X/Y}^T$, sont canoniquement iso entre eux.

( Il faut bien faire attention que $\mathcal E_{X/Y}$ \inserted{($Z$ non précisé)} est une classe d'iso. de groupes profinis sans plus, \uline{pas} un groupe extérieur. L'arbitraire dans le choix d'un iso entre $\mathcal E_{X/Y}^T$ et $\mathcal E_{X/Y}^{T'}$ provient du choix

\margin{112}
